\diarytitle{0123}{パラメータ表示された曲線に沿う線積分}

ベクトル場 $\bm{F} = (P, Q)$ に対する線積分は $\displaystyle\int_{C} \bm{F}\cdot d\bm{r} = \int_{C} P\,dx + Q\,dy$ で定義される。曲線 $C$ がパラメータ表示 $x = x(t)$, $y = y(t)$ $(a \le t \le b)$ で与えられるとき、$dx = x'(t)\,dt$, $dy = y'(t)\,dt$ を代入して $t$ の定積分に帰着させて計算せよ。

ベクトル場 $\bm{F}(x,y) = (y, x^{2})$ について、曲線 $C$: $x = t$, $y = t^{2}$ $(0 \le t \le 1)$ に沿う線積分
\[
\int_{C} \bm{F}\cdot d\bm{r} = \int_{C} y\,dx + x^{2}\,dy
\]
を求めよ。
